\paragraph{Б.1. Множества}
\subparagraph{Б.1.1}
Лень, картинки надо вставлять.

\subparagraph{Б.1.2}

Докажите обощение законов де Моргана для любого конечно набора множеств:

\[
\overline{A_1 \cap A_2 \cap \ldots A_n } = \overline{A_1} \cup \overline{A_2} \cup \ldots \overline{A_n }
\]
\[
\overline{A_1 \cup A_2 \cup \ldots A_n } = \overline{A_1} \cap \overline{A_2} \cap \ldots \overline{A_n }
\]

\textbf{Ответ:}



\subparagraph{Б.1.4}

Покажите, что множество нечетных натуральных чисел натуральных чисел счетно.

\textbf{Ответ:}

Определим множество все нечётных числе:
\[
X = \{ x : x \in N  \text{ и } x \bmod 2 = 1 \}
\]

Зададим биекцию $f : N \rightarrow X$:
\[
f(n) = 2n - 1
\]
Таким образом множество $X$ счётно.

\subparagraph{Б.1.5}

Покажите, что для любого конечного множества $S$ степенное множество $2^S$ содержит $2^{\left\| S \right\|}$ элементов (т.е. имеется $2^{\left\| S \right\|}$ различных подмножеств множества $S$).

\textbf{Ответ:}

По определению $2^S$ содержит все множества включая пустое множество и само множество $S$. Допустим $\left\| S \right\| = N$, тогда:
Общее кол-во множеств можно выразить:\\
0. Пустое множество - $\binom{N}{0}$ \\
1. Одноэлементных - $N$ или $\binom{N}{1}$; \\
2. Двух элементных - $\binom{N}{2}$; \\
3. Трех элементных - $\binom{N}{3}$; \\
и тд. \\
Общая сумма: \\
\[
2^{\left\| S \right\|} = \sum_{i = 0}^{N} {\binom{N}{i}} = 2^N = 2^{\left\| S \right\|}
\]

\subparagraph{Б.1.6}

Дайте индуктивное определение кортежа из $n$ элементов, используя понятие упорядоченное множество.

\textbf{Ответ:}

Рассмотрим базу индукции:
\[
(a, b) = \{\{a\}, \{a, b\}\}
\]

Выразим данное выражение как:

\[
(a, b) = ((a), b) = \{\{a\}, \{a, b\}\}
\]

Совершим индуктивный переход, тогда:

\[
(a_1, a_2, \cdots a_n, a_{n+1}) = ((a_1, a_2, \cdots a_n), a_{n+1})
\]
